\section{Unified Mobile and Shared-Television Continuity}
\label{sec:pilot-mobile-shared-tv}

The unified phone now acts as the private possession key for room-scale Pilot
surfaces.  The television remains a shared and comparatively untrusted display: it
does not receive phone private keys, provider credentials, private contact routes or
general authority merely because it is on the same network.

\subsection{Signed surface manifests}

The mother issues a short-lived Ed25519-signed home manifest for each exact surface.
The manifest identifies the surface, active persona or public state, issue and expiry
times, permitted tiles and privacy class.  Public mode contains only shared
experiences.  Tasks, calendar, files, personal work and the separately authenticated
Boardroom appear only while a trusted phone has acquired the television for that
persona.

Workspace tiles cannot be supplied as arbitrary interface data.  They are admitted
only from a separately valid provider-independent Workspace Manifest whose persona,
surface, membership capability and expiry match the television session.  Altered
labels, identities, capabilities or expiry values invalidate the signature.

The dashboard launcher is generated from this signed manifest rather than maintaining
a second independent set of private tiles.  The Devices experience remains the route
to topology, discovery and capability information; it is not duplicated as a global
technical navigation tab.

\subsection{Takeover, activity and logout}

Acquisition requires a nested possession signature from the same trusted phone that
holds the current mobile capability lease.  The session is exclusive: another adult
cannot silently replace the active person.  The television displays the active name,
and each accepted interaction renews a maximum five-minute inactivity window.
Explicit phone logout and the existing spoken logout route release the identity.

When a locked surface is rendered, its mother-side private presentation records and
expired private manifests are purged before a public manifest is issued.  Deliberate
document, dashboard or briefing presentation records retain the content identifier
and hash but not the document bytes.  Therefore logout removes the surface cache
without deleting the person's authoritative file, task or business record.

\subsection{Control and verification}

The phone reuses the established LG webOS gateway, room router, verified-navigation
authority and camera observer.  Observation and mutation have separate scopes.
Directional, media and allowlisted application controls require the active exclusive
session; a delivered button is not described as a visually successful outcome unless
subsequent evidence verifies it.  Device discovery is observe-only and cannot enroll
or control a newly seen device.

Sensitive message bodies and approval details remain private-phone data by default.
A manifest contains no such content.  A user must deliberately select a bounded
presentation, and the television receives no independent execution permission from
displaying it.

\subsection{Current qualification boundary}

The coded foundation covers the LG adapter, signed manifests, phone takeover,
identity display, inactivity expiry, logout, privacy filtering, cache purging, the
Devices route and room-aware command authority.  Full document rendering from every
Personal, Workspace and Boardroom repository remains a separate integration item, as
does representative multi-room field testing and non-LG hardware qualification.
Those are not inferred from unit tests.

Automated verification covers public and private tile separation, manifest
tampering, exact surface binding, deliberate presentation hashes, logout cache purge,
exclusive identity acquisition and the existing mobile TV authority.  Tests do not
operate the live household television.
